\section*{The Propers: Detailed Commentary}

\subsection*{Kingship, adoption, light, and truth\enspace\elementref{Entrance, Collect, Acclamation}}

Psalm 47's worldwide summons gives the Entrance a public horizon: the God acclaimed by all peoples is king over all the earth. The Collect narrows that grandeur to filial gift. Adoption is by grace, so “children of light” names received identity rather than spiritual pedigree. The petition to avoid error and adhere to truth also prevents light language from becoming mere inward feeling.

Augustine's \emph{Exposition on Psalm 47} receives the psalm's universal acclamation in relation to Christ and the gathered Church. This later ecclesial reading extends the psalm's Israelite worship toward the nations: the peoples themselves are summoned beneath the one kingship. The Collect then asks that a people adopted by this King live according to truth. Worldwide praise and filial obedience are thus related as gifts of the same King.

First Peter's acclamation receives Israel's royal-priestly vocabulary within a Christological address to baptized communities. Its purpose is proclamation of God's mighty acts, not self-congratulation.

The history of reception broadens that point. Clement of Alexandria's surviving comments on 1~Peter 2:9 distinguish election, royal belonging to Christ, priestly prayer, and the work of instruction. The \emph{Apostolic Constitutions} III.2 uses royal-priestly language in a baptismal and anointing setting. Augustine, in \emph{City of God} XVII, receives the Christian people as a royal priesthood under Christ the one priest, while Gregory the Great's \emph{Pastoral Rule} II uses the phrase to name the inward nobility and responsibility of regeneration. Clement gives the closest ancient comment on the verse; the others are illuminating ecclesial and pastoral reuse. None makes the acclamation a warrant for domination. Its grammar is vocational: those called from darkness are to announce another's mighty acts.

The witnesses also keep distinct what later Christian speech can too easily collapse. “Chosen” names divine initiative, “royal” derives from belonging to Christ, “priestly” orders prayer and witness, and proclamation names the end for which the people are gathered. Clement's sequence gives no basis for treating every baptized person as holding the same ecclesial office; neither does it permit reducing baptismal priesthood to a decorative title. Gregory's pastoral application makes privilege answerable to conduct: nobility of regeneration increases the obligation to govern oneself in virtue. Thus the acclamation prepares the Gospel not by declaring the hearers superior to those outside, but by reminding them why received light must become public praise and service.

\subsection*{A room, a promise, and a received messenger\enspace\elementref{First Reading, Psalm, Gospel}}

The Shunammite's hospitality is concrete and sustained: food, recognition, household consultation, and a small furnished room. The appointed omissions move from her welcome to Elisha's question about recompense and the servant's observation that she has no son. The promise is gift, not payment produced by a technique. The larger narrative, in which the son later dies and is restored, intensifies the story's dependence on God beyond the excerpt.

The narrative's details make discernment part of hospitality. The woman recognizes Elisha as a holy man through repeated encounter; she consults her husband; the room contains only what continuing welcome requires. Her practice is generous but ordered, neither spectacle nor indiscriminate access. When Elisha asks what might be done for her, she first answers from the security of her own people. Gehazi, not the woman, identifies childlessness. The promise therefore does not appear as the price she named for service.

No sustained verse-by-verse patristic commentary on precisely 2~Kings 4:8--11, 14--16a appears in the corpora searched, but several direct early uses of the appointed details do. Jerome, in \emph{Against Jovinianus} I.26, invokes the Shunammite who repeatedly welcomed Elisha and supplied table, bread, lamp, and the rest as an Old Testament example of women materially sustaining religious teachers. His ascetical polemic is not a commentary on her whole story, yet it recognizes her agency and the concrete provision named in the reading.

Origen's material preserved in \emph{Philocalia} XXVI names Elisha's very small chamber, bed, and inexpensive lamp in an argument against identifying righteousness with external wealth and health. The prophet who kept the law possessed little; the Shunammite's modest provision therefore becomes evidence against a prosperity reading. The work transmitted as the pseudo-Clementine \emph{Two Epistles on Virginity}, II.15, likewise remembers that she supported Elisha while the prophet lodged separately, then joins this detail to her later distress. Its Clementine attribution is disputed and its ascetical application reflects its own setting. Still, both witnesses directly receive the selected hospitality details, and both bear out that ordered boundaries and modest means do not diminish welcome.

Origen's use bears particularly on the promised son. He cites the sparse room to show that fidelity cannot be measured by external abundance; the later gift therefore cannot be converted into a universal contract of material return. The Shunammite's welcome is good before any recompense is named, and its goodness does not depend on forecasting the narrative's miracle.

Augustine's \emph{On the Soul and Its Origin} III.7 discusses only the later revival scene at 4:34, rejecting the claim that Elisha's breath became the child's soul. That is illuminating reuse of the sequel, not exegesis of the appointed unit. The reception is thus mixed rather than blank: identifiable Greek and Latin reception exists for the room and material support, but none of the witnesses found in the corpora searched gives a continuous exposition of the whole selected narrative.

Psalm 89 praises mercy forever, confesses a covenant, blesses those who know the festal cry, and identifies God as strength and shield. Yet the full psalm later laments apparent collapse of the Davidic promise. The response is therefore praise with eschatological tension, not evidence that covenant life is free of contradiction or grief.

Augustine's \emph{Exposition on Psalm 89} follows the covenant language toward Christ and the Church while attending to the psalm's apparent contradiction between promise and humiliation. His reception does not solve the historical lament by ignoring it; it asks how divine fidelity is discerned when visible royal success fails. In this liturgy, that question matters. The Shunammite receives a promise, but the unappointed sequel includes bereavement before restoration. Psalmic mercy is sung “for ever” from within temporal instability, not from outside it.

The appointed response takes only the psalm's praise and blessedness strophes, while the whole psalm goes on to the later protest, which the proclamation omits. That asymmetry protects the promise to the Shunammite from becoming a rule that hospitality yields uninterrupted household flourishing. Augustine's attention to apparent covenant failure holds praise, lament, and eschatological trust together rather than letting one cancel the others.

Matthew ends the missionary discourse by making reception relationally deep: receiving the apostolic messenger receives Christ, and receiving Christ reaches the Father who sent him. Chrysostom guards the dignity of apparently small persons and acts. Prophet, righteous person, disciple, and giver of water all enter a chain in which Christ takes the treatment of his servants personally.

Chrysostom's \emph{Homily 35 on Matthew} treats both the messenger's dignity and the giver's poverty. Christ offers reward even for cold water so that lack of wealth cannot become an excuse for withholding charity. In \emph{Homily 31 on Hebrews}, Chrysostom later reuses the cup saying to make the same pastoral point about hospitality: small means do not abolish responsibility or the value Christ assigns to welcome. Cyprian's \emph{Testimonies} III.42 likewise includes the cup among works whose apparently small scale does not escape divine notice.

Two Latin receptions refine the act's end. Augustine's \emph{Exposition on Psalm 122}, section 10, insists that the value of cold water lies in charity rather than quantity. John Cassian's \emph{Conference 1}, chapter 10, uses the promised reward to distinguish temporary works from charity, the enduring goal toward which they serve. These are illuminating reuses rather than further expositions of Matthew's entire discourse. Together they keep reward from becoming commerce: the gift is small enough for the poor, significant because offered in discipleship, and ordered toward charity rather than spiritual accounting.

This does not confer unlimited access or make every claimant a trustworthy messenger. Ecclesial authority, discernment, safeguarding, and prudence remain necessary. Nor is reward commercial payment. The cup matters because it is given in the name of discipleship; grace dignifies the person and act beyond their visible scale.

\subsection*{Christ before family and the cross without self-harm\enspace\elementref{Gospel}}

Jesus' comparative saying demands supreme allegiance. Chrysostom interprets it as preferring the Creator when creaturely attachment conflicts with discipleship, not hating parents as such. Elsewhere Matthew condemns evasion of duties to father and mother. The hard saying therefore orders loves: family is loved truthfully under God, never made an absolute rival and never discarded as though grace opposed nature.

Taking one's cross is readiness for faithful suffering in union with Christ. It does not sacralize abuse, disable self-protection, or turn preventable injury into a spiritual achievement. Losing life for Christ can include martyrial fidelity and daily self-gift, but cannot be invoked to command another person's silence or exposure.

\subsection*{Baptized into death and life\enspace\elementref{Second Reading}}

Romans 6 begins from sacramental fact. Baptism joins believers to Christ's death and burial so that their life may correspond to his resurrection. Paul is not describing a private symbol chosen after conversion; he appeals to a common ecclesial initiation that changes the believer's relation to sin and to Christ.

Chrysostom's \emph{Homilies 10--11 on Romans} follows Paul's rebuttal of the idea that abundant grace makes continued sin harmless. Baptismal burial is not a theatrical resemblance: Christ died bodily once, while the baptized are to die to sin and live in a new manner. Aquinas's \emph{Commentary on Romans} VI, lectio 1, similarly distinguishes sacramental sign, incorporation into Christ's death, and the moral consequence of newness of life. Both witnesses treat the passage directly. Their common reading is prospective as well as retrospective: baptism names what God has done and therefore the form life must now take.

The return at verses 8--11 adds an eschatological confidence that moral exhortation alone could not supply. Believers trust that they will live with Christ because death no longer has dominion over the risen Christ; their reckoning of themselves dead to sin rests on his once-for-all act. “Consider yourselves” is therefore neither make-believe nor instant moral perfection. It is judgment according to sacramental reality, repeatedly enacted against sin's attempted return. That logic can illuminate Matthew's language of losing and finding life, but only as a later canonical synthesis: Romans remains the independent semi-continuous reading.

The appointed boundary has a gap. Verses 5--7, with further explanation of union and freedom, are omitted; verses 8--11 resume the movement to living with Christ and reckoning oneself dead to sin. Chrysostom distinguishes Christ's bodily death from the baptized person's death to sin, while joining both through sacramental participation. The once-for-all death of Christ grounds rather than removes ongoing fidelity.

\Needspace{5\baselineskip}
Romans remains semi-continuous. A cross--Baptism connection with Matthew is theologically plausible and canonically grounded.

\subsection*{Mysteries effected by God\enspace\elementref{Prayer over Offerings}}

The Prayer over the Offerings assigns sacramental effect to God and asks that the Church's service accord with the mysteries. Divine efficacy and human disposition are related without competition: God acts first and truly; worshippers are genuinely called to worthy service.

\subsection*{Remembered benefits, unity, life, and charity\enspace\elementref{Communion, Prayer after Communion}}

Communion option A begins Psalm 103's self-address: the whole person blesses God and refuses forgetfulness of his benefits. Option B carries Christ's prayer that later believers be one so the world may believe. Both are appointed textual alternatives.

Augustine's \emph{Exposition on Psalm 103} receives the command to one's own soul as disciplined memory: praise gathers the person by recalling divine benefits rather than waiting for an unbroken feeling of gratitude. The first branch thus complements the Prayer after Communion's petition for continuing charity, but only if selected.

Option B has unusually close direct Johannine reception. Augustine's \emph{Tractate 109 on John} emphasizes that Jesus prays not only for the apostles present but for future believers who come through their word. \emph{Tractate 110} then distinguishes the unity requested for creatures from the Father's and Son's unity of substance: ecclesial communion is a graced participation and likeness, not fusion or erased difference. The missionary clause---that the world may believe---makes unity outward-facing, but it never makes appearance more important than truth. These tractates illuminate the branch; its use on the occurrence is unresolved.

The Prayer after Communion names the Eucharistic gift as the offered and received divine victim, asking life and lasting charity. Option B sharpens ecclesial unity; option A, grateful memory. The prayer governs both.

\section*{The Propers: Notable and Quotable}

{\small
\begin{itemize}
\item \textbf{“Take up his cross” --- from discipleship to lawn-mowing.\enspace\elementref{Gospel}} The English idiom “a cross to bear” now names any stubborn burden or trial. It can remain grave, but Ammer's \emph{American Heritage Dictionary of Idioms} also illustrates deliberately light uses. A saying about allegiance unto suffering has become elastic enough to dignify an oversized weekly lawn.

\item \textbf{“Let us arrange a little room on the roof” --- the Prophet's Chamber.\enspace\elementref{First Reading}} At Trewint in Cornwall, Digory Isbell read the Shunammite's story and resolved to imitate it. He added two rooms for John Wesley and itinerant preachers; the upper room held the biblical bed, table, stool, and candlestick. A prophet's lodging became an architectural institution for the next travelling preacher.

\item \textbf{“A cup of cold water” --- Edith Wharton's title without the water.\enspace\elementref{Gospel}} Wharton's 1899 story gives the phrase to a compromised bank clerk who interrupts a stranger's suicide, pays her lodging and train fare, and then abandons his own flight to face those he has defrauded. No literal cup appears. The tiny Gospel gift becomes a title for one mercy that redirects two lives, while the giver's guilt keeps the moral reversal sharply ironic.
\end{itemize}
}

\section*{The Propers: Interpretive Possibilities}

\noindent\emph{The following exploratory proposals were developed in an AI-assisted editorial process. They are attributed to none of the cited authorities and claim no historical compositional intent. Each states its anchors and limits.}

\Needspace{8\baselineskip}
\noindent\textbf{Hospitality becomes durable when charity enters architecture.\enspace\elementref{First Reading; Gospel}}
The 2~Kings narrative turns recognition into a repeatable furnished room; the Gospel places reception within anyone's reach as cold water. Their mechanism is configured material charity---space, time, budget, access, and safeguards that sustain another's mission. Read separately, generous household and small gift miss ordinary holiness's design intelligence. The fruit is an audit of what recurring good homes and parishes have made easy. Ancient witnesses preserve modest room and accessible cup, but no checked source gives this infrastructure account. Welcome is not indiscriminate access, reward is not commerce, and the promised child is no technique or guarantee.

\Needspace{8\baselineskip}
\noindent\textbf{Royal priesthood can appear as a lay hand holding water.\enspace\elementref{Entrance; Acclamation; Gospel}}
The Entrance's universal King, the acclamation's royal priesthood ordered to proclamation, and the Gospel's hand receiving a messenger with water form a descending scale. Its mechanism makes small material care a public enactment of derivative kingship and priestly mediation; separate reading leaves dignity abstract and hospitality private. Chrysostom supplies a striking precedent: the lay giver becomes a priest of Christ, yet first receives all from Christ. The fruit is baptismal agency without clerical confusion. This neither equates common and ministerial priesthood nor erases governance, competence, or safeguarding.
